\documentclass[reqno, answers]{exam}

\usepackage{amssymb, amsmath, amsthm, stmaryrd}
\usepackage{fullpage, parskip}
\usepackage{enumitem}
\usepackage{tikz}
\usepackage{ulem}
\usepackage{hyperref}

\title{Introduction to Computational Math Spring 2026 \\ Homework 07}
\date{Due: Friday, Mar 13, 2026, 11:59 PM}
\author{}

\begin{document}
\maketitle
\thispagestyle{empty}

Textbook = Driscoll and Braun (\url{https://fncbook.com/})

\begin{enumerate}

    \item Textbook Exercise 5.1.4
          \begin{solution}
              \textbf{(a)} We verify that $q(s)$ passes through the three points by direct substitution:

              At $s = -1$:
              $$q(-1) = a \cdot \frac{(-1)(-2)}{2} - b(-2)(0) + c \cdot \frac{(-1)(0)}{2} = a \cdot 1 - 0 + 0 = a$$

              At $s = 0$:
              $$q(0) = a \cdot \frac{0(-1)}{2} - b(-1)(1) + c \cdot \frac{0(1)}{2} = 0 + b + 0 = b$$

              At $s = 1$:
              $$q(1) = a \cdot \frac{1(0)}{2} - b(0)(2) + c \cdot \frac{1(2)}{2} = 0 - 0 + c = c$$

              Therefore, $q(s)$ interpolates the points $(-1,a)$, $(0,b)$, $(1,c)$.

              \textbf{(b)} We need $s = Ax + B$ such that:
              \begin{align*}
                  s = -1 & \text{ when } x = x_0 - h \\
                  s = 0  & \text{ when } x = x_0     \\
                  s = 1  & \text{ when } x = x_0 + h
              \end{align*}

              From $s = 0$ when $x = x_0$: $B = -Ax_0$.

              From $s = 1$ when $x = x_0 + h$: $1 = A(x_0 + h) + B = Ah$, so $A = \frac{1}{h}$.

              Therefore: $$s = \frac{x - x_0}{h}$$

              \textbf{(c)} Substituting $s = \frac{x - x_0}{h}$ into $q(s)$:
              \begin{align*}
                  \tilde{q}(x) & = a \cdot \frac{s(s-1)}{2} - b(s-1)(s+1) + c \cdot \frac{s(s+1)}{2}
              \end{align*}

              Expanding and simplifying with $s = \frac{x-x_0}{h}$:
              $$\tilde{q}(x) = \frac{a(x-x_0)(x-x_0-h) - 2b(x-x_0-h)(x-x_0+h) + c(x-x_0)(x-x_0+h)}{2h^2}$$

              Alternatively, in compact form:
              $$\tilde{q}(x) = \frac{(a-2b+c)(x-x_0)^2 + (c-a)h(x-x_0) + 2bh^2}{2h^2}$$
          \end{solution}

    \item Textbook Exercise 5.1.5
          \begin{solution}
              From part (a) of the previous exercise, we can write $q(s)$ in terms of cardinal functions:
              $$q(s) = a\phi_{-1}(s) + b\phi_0(s) + c\phi_1(s)$$
              where
              \begin{align*}
                  \phi_{-1}(s) & = \frac{s(s-1)}{2}, \quad \phi_0(s) = -(s-1)(s+1), \quad \phi_1(s) = \frac{s(s+1)}{2}
              \end{align*}

              We compute the infinity norms on $s \in [-1, 1]$:

              For $\phi_{-1}(s) = \frac{s(s-1)}{2}$: This is non-positive on $[-1,1]$ with minimum at $s=1/2$ giving $-1/8$. Thus $\|\phi_{-1}\|_\infty = 1/8$.

              For $\phi_0(s) = -(s-1)(s+1) = 1 - s^2$: Maximum at $s=0$ gives $1$. Thus $\|\phi_0\|_\infty = 1$.

              For $\phi_1(s) = \frac{s(s+1)}{2}$: This is non-negative on $[-1,1]$ with maximum at $s=1$ giving $1$, but checking $s=-1$ gives $0$ and the critical point $s=-1/2$ gives $-1/8$. Actually, on $[-1,1]$, the maximum magnitude is at $s=1$ giving $1$. Thus $\|\phi_1\|_\infty = 1$.

              By Theorem 5.1.1, the condition number satisfies:
              $$\max\left\{\frac{1}{8}, 1, 1\right\} \leq \kappa(y) \leq \frac{1}{8} + 1 + 1$$

              Therefore:
              $$1 \leq \kappa(y) \leq \frac{17}{8}$$

              The bounds show that quadratic interpolation at equally-spaced nodes $\{-1, 0, 1\}$ is well-conditioned, with the condition number bounded by $17/8 = 2.125$.

              By the change of variable in part (b), these bounds apply equally to interpolation at nodes $\{x_0-h, x_0, x_0+h\}$ for any $x_0$ and $h > 0$.
          \end{solution}

    \item Textbook Exercise 5.2.5

                    \begin{solution}
              \textbf{(a)} Apply the mean value theorem to $f$ on the interval $[t_k, x] \subset (t_k, t_{k+1})$:
              $$f'(s) = \frac{f(x) - f(t_k)}{x - t_k}$$
              for some $s \in (t_k, x) \subset (t_k, t_{k+1})$.

              Since $f(t_k) = y_k$, we have:
              $$f(x) = y_k + (x - t_k)f'(s)$$

              \textbf{(b)} The piecewise linear interpolant satisfies:
              $$p_n(x) = y_k + (x - t_k)\frac{y_{k+1} - y_k}{t_{k+1} - t_k}$$

              Apply the mean value theorem to $f$ on $[t_k, t_{k+1}]$:
              $$f(t_{k+1}) - f(t_k) = (t_{k+1} - t_k)f'(u)$$
              for some $u \in (t_k, t_{k+1})$, so:
              $$f'(u) = \frac{y_{k+1} - y_k}{t_{k+1} - t_k}$$

              Now apply the mean value theorem to $f'$ on the interval with endpoints $s$ and $u$:
              $$f''(v) = \frac{f'(s) - f'(u)}{s - u}$$
              for some $v$ between $s$ and $u$, hence $v \in (t_k, t_{k+1})$.

              This gives:
              $$f'(s) - \frac{y_{k+1} - y_k}{t_{k+1} - t_k} = (s - u)f''(v)$$

              \textbf{(c)} From parts (a) and (b):
              \begin{align*}
                  f(x) - p_n(x) & = y_k + (x - t_k)f'(s) - y_k - (x - t_k)\frac{y_{k+1} - y_k}{t_{k+1} - t_k} \\
                                & = (x - t_k)\left[f'(s) - \frac{y_{k+1} - y_k}{t_{k+1} - t_k}\right]         \\
                                & = (x - t_k)(s - u)f''(v)
              \end{align*}

              Taking absolute values and using $|f''(v)| \leq M = \|f''\|_\infty$:
              $$|f(x) - p_n(x)| \leq |x - t_k| \cdot |s - u| \cdot M$$

              Since $x, s, u \in (t_k, t_{k+1})$, we have $|x - t_k| < h$ and $|s - u| < h$.

              Therefore:
              $$\|f - p_n\|_\infty \leq Mh^2$$
          \end{solution}

    \item Jupyter notebook

\end{enumerate}


\section*{Quiz Information}

The quiz will be held on {\bf Tuesday, March 24, 2026} during the discussion section.

Quiz questions will be modifications of the homework problems and material discussed in class.

\textbf{Topics covered:}

Chapters 5.1 and 5.2.

\begin{itemize}
    \item Conditioning of linear interpolants
    \item Cardinal functions
    \item Properties of linear interpolation
\end{itemize}

The quiz on this material will be short as we covered very little material.



\end{document}